\begin{solapartat}
\punts{0,1} La longitud d'ona és la distància mínima que separa dues crestes: $\lambda = 4\ \mathrm{m}$.

\punts{0,1} La freqüència és el nombre de cicles que es produeixen per segon:
\[ f = \tfrac{30\ \mathrm{cicles}}{60\ \mathrm{s}} = 0,50\ \mathrm{Hz}. \]

\punts{0,2} La velocitat de propagació de l'ona és: $v = \lambda f = 2,00\ \mathrm{m/s}$.

\punts{0,2} La boia descriurà un MHS:
\[ y(t) = A\sin(\omega t + \varphi_0) \]

\punts{0,1} La separació entre una cresta i una vall és dues vegades l'amplitud:
\[ A = 0,20\ \mathrm{m}. \]

\punts{0,1} $\omega = 2\pi f = \pi\ \mathrm{rad/s}$

\punts{0,2}
\[ A = A\sin(\varphi_0) \Rightarrow \varphi_0 = \mathrm{ArcSin}(1) = \pi/2\ \mathrm{rad}. \]

\punts{0,25} Així, l'equació del moviment de la boia és:

$y(t) = 0,2\sin\left(\pi t + \frac{\pi}{2}\right)$, $t$ en s, $y$ en m.

Alternativament, es pot descriure l'MHS així: $x(t) = A\cos(\omega t + \varphi_0)$. En aquest cas, $\varphi_0 = 0\ \mathrm{rad}$.
\end{solapartat}

\begin{solapartat}
\punts{0,5} $v(t) = \frac{dy(t)}{dt} = A\omega\cos(\omega t + \varphi_0)$,

Llavors:

$v(t) = 0,628\ \cos\left(\pi t + \frac{\pi}{2}\right)$, t en s, v en m/s.

Atès que el valor màxim de cosinus és 1:
\[ v_\text{màx} = A\omega = 0,628\ \mathrm{m/s} \]

\punts{0,25}
\begin{gather*}
a(t) = \frac{dv(t)}{dt} = -A\omega^2\sin(\omega t + \varphi_0) = -y(t)\,\omega^2\\
a(t) = -1,97\sin\left(\pi t + \frac{\pi}{2}\right),\ \text{t en s, a en m/s}^2.
\end{gather*}
\begin{center}
Atès que el valor màxim de $-\sin(\omega t + \varphi_0)$ és 1:
\end{center}
\[ a_\text{màx} = A\omega^2 = 1,97\ \mathrm{m/s^2} \]

\punts{0,5} \destaca{Representació gràfica:}
\figura[0.5]{sol_fig1.png}
No incloure el títol en l'eix resta \destaca{0,1 p.}

No incloure les unitats en l'eix resta \destaca{0,2 p.}

Si, en la representació, la fase inicial no es correspon amb la de l'equació deduïda resta \destaca{0,2 p.}

No escalar o representar correctament la gràfica serà qualificat amb un zero.
\end{solapartat}
